We describe an extension of \galpy, the \python\ library for galactic
dynamics, that adds \jax\ and \torch\ as computational backends
alongside its existing \numpy\ and \cc\ implementation. With it,
\galpy's gravitational potentials, orbit integration, action--angle
transformations and their inverses, and distribution functions,
including those for tidal streams, run on GPUs, can be compiled as part
of a user's program, and can be differentiated with respect to their
inputs and parameters, from the force at a point to the samples drawn
from a distribution function. The user-facing design is kept simple: the
active backend follows from the type of the arrays that are passed in,
the \galpy\ interface is unchanged, and the \numpy\ path gives the same
results as before. We describe the design of the
backend framework and its implementation in each of \galpy's modules,
with examples of use, and we discuss its limitations. On a GPU, batches
of orbits, samples from distribution functions, and stream models run
one to five orders of magnitude faster than on a single CPU core. To
illustrate the new capabilities, we present simple examples of fitting
the Galactic potential with a hypervelocity star and of learning the
inverse action--angle transformation with a neural network trained
through \galpy. An AI coding agent wrote
essentially all of the code under the detailed direction and review of
the human author, and we describe how we organized this work and how we
established the correctness of the result. This work is the start of
\galpy's version~2, which is available at
\url{https://github.com/jobovy/galpy} and documented at
\url{https://docs.galpy.org/en/v2/}.
